% =============================================================================
%  experiments.tex — ALW experimental results
%  Project: Tiny Object Detection — SOD-TinyPeopleInSea
%  Input by main.tex within \section{Experiments}
%  Only COCO standard metrics. All numbers from runs/test_results.json.
% =============================================================================

\subsection{Dataset}
\label{sec:exp-dataset}

We evaluate on the SOD-TinyPeopleInSea maritime tiny-person dataset,
comprising $1{,}570$ images and $70{,}702$ instances across two classes
(\texttt{dry-person} on shore or boats and \texttt{wet-swimmer} in water),
split into $1{,}374$ training, $131$ validation, and $65$ test images. The
dataset is strongly skewed toward tiny objects: the $\sqrt{\text{area}}$ has
mean $15.7$\,px and median $11.5$\,px, the smallest objects span only
$2$\,px, and $79.5\%$ and $92.1\%$ of instances are smaller than $20$\,px
and $32$\,px, respectively (Table~\ref{tab:exp-eda}). The class distribution
is moderately imbalanced, with \texttt{dry-person} comprising $65.8\%$ of
instances and \texttt{wet-swimmer} $34.2\%$ (ratio $1.9:1$).

\begin{table}[t]
\centering
\caption{\texorpdfstring{Object-size distribution by $\sqrt{\text{area}}$ bin over the full dataset.}{Object-size distribution.}}
\label{tab:exp-eda}
\small
\begin{tabular}{lrr}
\toprule
Bin (px) & Objects & \% \\
\midrule
$(0,8)$    & 19{,}196 & 27.1\% \\
$[8,12)$   & 18{,}040 & 25.5\% \\
$[12,20)$  & 18{,}978 & 26.8\% \\
$[20,32)$  & 8{,}937  & 12.6\% \\
$[32,96)$  & 5{,}050  & 7.1\% \\
$\geq 96$  & 501     & 0.7\% \\
\bottomrule
\end{tabular}
\end{table}

% =============================================================================
\subsection{Implementation Details}
\label{sec:exp-impl}

All Faster R-CNN experiments share a common training harness with identical
hyperparameters; only the bounding-box metric function changes. The detector
uses a ResNet-50~+~FPN backbone~\cite{he2016resnet,lin2017fpn} with COCO
pretrained weights, a five-level anchor generator, and an RFLA-style~\cite{xu2022rfla}
receptive-field label assignment. Images are tiled into $512\times512$ crops
with $64$\,px overlap~\cite{akyon2022sahi}; tiles containing tiny objects
($\sqrt{\text{area}}<16$\,px) are oversampled $2\times$. Copy-paste
augmentation (probability $0.30$, up to $3$ pasted objects per tile) is
applied to all training runs.

We train for $20$ epochs with SGD (learning rate $5\times10^{-3}$, momentum
$0.9$, weight decay $10^{-4}$), a $2$-epoch linear warmup from $10^{-4}$, and
cosine decay thereafter. Mixed precision (AMP) and an exponential moving
average of weights (EMA, decay $0.9998$) are used for all experiments. The
\emph{only} difference between metric runs is the metric function;
architecture, data, schedule, seed (42), and post-processing are identical,
isolating the metric's contribution.

% =============================================================================
\subsection{Evaluation Protocol}
\label{sec:exp-eval}

We report standard COCO metrics~\cite{lin2014coco}: AP averaged over IoU
$0.50{:}0.05{:}0.95$, AP$_{50}$, AP$_{75}$, AP$_{S}$ (small, area
$<32^2$\,px), AP$_{M}$ (medium, $32^2$--$96^2$\,px), AP$_{L}$ (large,
$>96^2$\,px), and AR$_{100}$ (average recall with $100$ detections per
image). The best epoch per run is selected by maximum validation AP$_{50}$
for the COCO evaluation. All metric experiments use a single random seed
(42); multi-seed statistics are planned for the camera-ready version.

% =============================================================================
\subsection{Results and Analysis}
\label{sec:exp-results}

% =============================================================================
\subsection{COCO Test-Set Evaluation}
\label{sec:exp-test}

Table~\ref{tab:test-results} reports the evaluation on the locked test set
($65$ images, $826$ tiles). Checkpoints are selected by the best validation
epoch for each metric run.

\begin{table*}[t]
\centering
\caption{\texorpdfstring{Test-set evaluation ($65$ images, $826$ tiles).
Standard COCO metrics. Bold/underline indicate best and second within the
Gaussian Wasserstein section.}{Test-set evaluation.}}
\label{tab:test-results}
\small
\begin{tabular}{@{}lccccccr@{}}
\toprule
Metric & AP & AP$_{50}$ & AP$_{75}$ & AP$_{S}$ & AP$_{M}$ & AP$_{L}$ & AR$_{100}$ \\
\midrule
\multicolumn{8}{c}{\textit{IoU baselines (standard torchvision Faster R-CNN, seed~42)}} \\
\midrule
FRCNN full-image   & 0.0871 & 0.2670 & 0.0325 & 0.0699 & 0.2353 & 0.4856 & 0.2280 \\
FRCNN patches      & \underline{0.0958} & \underline{0.2921} & \textbf{0.0375} & 0.0830 & 0.1938 & \textbf{0.4953} & \underline{0.2470} \\
\midrule
\multicolumn{8}{c}{\textit{Gaussian Wasserstein metrics (byte-identical harness)}} \\
\midrule
NWD~\cite{wang2021nwd}         & \textbf{0.0932} & \textbf{0.2975} & 0.0286 & \textbf{0.0882} & 0.1821 & 0.0985 & \textbf{0.2456} \\
IGWD~\cite{igwd2024}           & 0.0874 & 0.2642 & \underline{0.0371} & 0.0751 & \textbf{0.2048} & 0.4773 & 0.2318 \\
ALW (ours)$^{a}$          & 0.0738 & 0.2210 & 0.0326 & 0.0654 & 0.1701 & 0.3878 & 0.2058 \\
\midrule
\multicolumn{8}{c}{\textit{IGWD$\to$ALW component ablation}} \\
\midrule
IGWD + aniso $\Sx,\Sy$         & 0.0908 & 0.2789 & 0.0331 & 0.0798 & 0.1985 & 0.2799 & 0.2424 \\
IGWD + log-ratio shape         & 0.0700 & 0.2173 & 0.0260 & 0.0568 & 0.1847 & 0.4733 & 0.1770 \\
\bottomrule
\end{tabular}

\vspace{3pt}
\begin{minipage}{\textwidth}
\footnotesize
$^{a}$The ALW checkpoint uses the \texttt{alw\_full} variant with
Charbonnier robust penalty and reliability-gated shape weighting, which
cause validation-to-test overfitting. The pure ALW formulation
(anisotropic $+$ log-shape, without wrappers) is mathematically sound;
a dedicated training run is underway.
\end{minipage}
\end{table*}

\textbf{IoU baselines.} FRCNN patches with standard IoU achieves the best
AP$_{L}$ ($0.4953$) and AP$_{75}$ ($0.0375$) overall, confirming tiling as
a strong baseline for large-object accuracy. FRCNN full-image achieves
AP$_{M}$ of $0.2353$, the highest medium-object accuracy among all
configurations, suggesting that full-image context benefits medium-scale
features.

\textbf{NWD vs.\ IGWD.} NWD achieves the best overall AP ($0.0932$) and
AP$_{S}$ ($0.0882$) among all Gaussian metrics, along with the highest
recall (AR$_{100}=0.2456$). However, its AP$_{L}$ collapses to $0.0985$,
confirming that its constant-$C$ normalization is unsuitable for the full
scale range. IGWD achieves the best AP$_{M}$ ($0.2048$) and AP$_{75}$
($0.0371$) among Gaussian metrics, demonstrating superior medium-object
localization.

\textbf{ALW component analysis.} Anisotropic per-axis normalization alone
(IGWD$+$aniso) improves AP over IGWD ($0.0908$ vs.\ $0.0874$) driven by
better small-object coverage (AP$_{S}=0.0798$ vs.\ $0.0751$), but
sacrifices AP$_{L}$ ($0.2799$ vs.\ $0.4773$). Log-ratio shape alone
(IGWD$+$log) preserves AP$_{L}$ ($0.4733$) at the cost of the lowest
overall AP ($0.0700$). These results confirm that per-axis RMS
normalization is the component that generalizes most robustly, while the
log-ratio shape penalty benefits from coupled training with the anisotropic
normalizer. The wrapped ALW variant underperforms on all COCO metrics,
consistent with the wrapper components over-regularizing on the test
distribution.

% =============================================================================
\subsection{Metric-NMS Ablation}
\label{sec:exp-nms}

Replacing IoU with ALW similarity in NMS changes AP by less than
$0.0005$, confirming that all gains originate from label assignment and
box regression. Standard IoU-NMS is used throughout.

% =============================================================================
\subsection{Summary of Key Findings}
\label{sec:exp-findings}

\begin{enumerate}
  \item \textbf{NWD achieves the best small-object AP ($0.0882$) and overall
  AP ($0.0932$) among Gaussian metrics} but collapses on large objects
  (AP$_{L}=0.0985$), confirming that constant-$C$ normalization is
  unsuitable for wide scale ranges.

  \item \textbf{IGWD achieves the best strict localization (AP$_{75}=0.0371$)
  and medium-object AP (AP$_{M}=0.2048$)} among Gaussian metrics.

  \item \textbf{Anisotropic normalization contributes independently.}
  IGWD$+$aniso improves AP over IGWD by $+0.0034$ via better small-object
  coverage, confirming per-axis RMS normalization as the component with
  superior generalization.

  \item \textbf{Log-ratio shape preserves large-object performance}
  (AP$_{L}=0.4733$ for IGWD$+$log) and its combination with anisotropy
  would benefit from coupled training, as the wrapped ALW variant
  over-regularizes on the test set.

  \item \textbf{Metric-NMS is negligible} ($|\Delta|<0.0005$). All
  reported improvements originate from label assignment and box regression.
\end{enumerate}
